\section{Methods}
\label{sec:methods}

\subsection{Problem and required output}

All backends solve the same linear program: maximise $\sum_i p_i x_i$ subject to
$\sum_i w_i x_i \le c$, $x_i \le x_j$ for every arc $(i,j)$ meaning ``$i$
requires $j$'', and $0 \le x_i \le 1$. Profits carry either sign, weights are
strictly positive, and there is a single capacity row. Cyclic digraphs are
accepted and condensed; solutions and multipliers are expanded back to the
original variables and arcs.

The required output of a \texttt{solve} task is a primal solution, the
multipliers $\lambda$, $\mu$ and $\alpha$, and the canonical blocks of the
positive path. Comparisons are therefore between runs that produce the same
object: whenever a configuration answers several capacities from one path, the
report says so explicitly and section~\ref{sec:scalability} treats it separately.

\subsection{Compared methods}

A note on the labels. The method developed by this project is the \FS{}, and
the configuration labels in the current suites are named after it,
\texttt{ratio-closure-*}. The labels below instead read \texttt{forest-*},
because that is what the campaigns reported here actually recorded: they were
measured on 6 September 2026, before decision~D18 was applied to the code, and a
run identifier is derived from the instance, the task, the capacity, the
configuration \emph{content} and the hash of the sources. The manifests are a
record of what was executed and are not rewritten, so this report quotes the
labels as measured. The divergence is cosmetic and disappears with the first
campaign measured under the current suites; the analysis script recognises both
spellings.

\begin{description}[leftmargin=2.4em,style=nextline]
  \item[\texttt{forest-reference}] \FS{} with full pricing, Bland's entering
    rule, complete forest rebuilds and sequential canonicalisation. This is the
    verifiable reference implementation, not a configuration meant to be fast.
  \item[\texttt{forest-tuned}] the configuration frozen by the random search of
    section~\ref{sec:tuning}, chosen on training and validation data only.
  \item[\texttt{dinkelbach}] Newton iteration on
    $\max\{p(S)-\lambda w(S)\}$ over closures, started from a strict lower bound
    on the ratio so that negative ratios need no special case. Each iterate is a
    maximum-closure minimum cut.
  \item[\texttt{parametric}] divide and conquer on the interval of ratios: the
    mean ratio of an interval is either a breakpoint or splits it in two. Each
    subproblem is an independent maximum flow.
\end{description}

The maximum-flow routine is a Dinic implementation written for this project
(level graphs, current arc, iterative augmentation). It is correct and
reasonable, but it is \emph{not} a state-of-the-art flow code: no pseudoflow, no
highest-label push-relabel with gap and global relabelling, no reuse of flow
across parametric subproblems. Any statement about the flow backends in this
report is a statement about this implementation, and
section~\ref{sec:discussion} says what that does and does not license.

\subsection{Accuracy is fixed before speed}

All configurations use the same tolerances, and tolerances are not part of the
tuning space. Every run must pass, in the original units, the primal checks
($0 \le x \le 1$, every original arc, the capacity row), the dual sign
conditions, the reduced-cost inequality
$w_i\lambda + \mu_i + \divg_i(\alpha) \ge p_i$, complementarity on $\mu$, on
$\alpha$ and on the capacity row, and equality between the primal objective and
the dual bound. A result that fails any of these is reported as
\texttt{numerical\_failure} and never counted as a success.

Verification is a separate process: after each run, \texttt{pclp verify} reloads
the instance and the saved result and re-checks the certificate without
re-optimising. Only runs that pass this second, independent check enter the
analysis.
